
\usepackage[frenchb]{babel}
\usepackage[utf8]{inputenc}
\usepackage[T1]{fontenc}
\usepackage{lipsum}
\usepackage{graphicx}
\usepackage[left=3.5cm, right=2cm, top=2cm, bottom=2.5cm]{geometry}
\usepackage{listings}
\usepackage[hidelinks]{hyperref}
\usepackage{amsmath,amsfonts,amssymb}
\usepackage{tikz}
\usepackage{shorttoc}
\usepackage{makeidx}
\usepackage[nopostdot,nonumberlist,toc]{glossaries}
\usepackage{nomencl}
\usepackage{tabularx}
 \usepackage{ifthen}
\usepackage[Lenny]{fncychap}

\author{Mathieu SOUM}

\makeatletter
\let\thetitle\@title
\let\theauthor\@author
\makeatother

\newcommand{\bigtitle}{Participation au développement d'une application Java EE}
\newcommand{\shorttitle}{dédiée à la conception et à la diffusion de supports pédagogiques}
\newcommand{\tuteur}{Sébastien RICHE}
\newcommand{\encadrant}{Isabelle FERRAN\'E}
\newcommand{\coverfoot}{Master 2 -- Développement Logiciel / 2014 - 2015}

\newcommand{\coverpage}[1]{%
\newgeometry{top=2cm,left=2cm,bottom=2cm,right=1cm}
\begin{titlepage}
\parindent=0pt
\begin{tabular}[h]{l@{ : }l}
  Tuteur & \tuteur \\
  Encadrant & \encadrant
\end{tabular}
\ifthenelse{\equal{#1}{cover}}%
{\vspace{1.5cm}}%
{\vspace{3.5cm}}
\begin{center}
\ifthenelse{\equal{#1}{cover}}
{
  \includegraphics[scale=0.9]{images/aboard.png}%
}{}
\end{center}
\vspace{2cm}
\hrulefill
\begin{center}\bfseries\Huge
  \bigtitle
\end{center}
\hrulefill
\begin{center}\bfseries\Large
  \shorttitle
\end{center}
\vspace*{0.5cm}
\begin{center}\bfseries\Large
  \theauthor
\end{center}
\vspace*{0.2cm}
\begin{center}\bfseries
  Stage du 3 mai 2014 au 31 juillet 2015\\Chez Aboard Engineering
\end{center}
\vfill
\ifthenelse{\equal{#1}{cover}}
{
\begin{tabular}[h]{m{5.7cm}m{5.7cm}m{5.7cm}}
\includegraphics[scale=0.45]{images/aboard.png} &
\includegraphics[scale=0.55]{images/ups.png} &
\includegraphics[scale=0.55]{images/mdl.png} \\
\end{tabular}
\vspace{1.2cm}
}{}
\begin{flushright}
  \coverfoot
\end{flushright}
\end{titlepage}
\restoregeometry
}

\makeindex

%\fancyhead[RE]{\rightmark}
%\fancyhead[LO]{\leftmark}
%\fancyhead[RO,LE]{}
%\fancyfoot[LE,RO]{\thepage}
%\fancyfoot[C]{}

\usepackage{fancyhdr}%pour les en-têtes et pieds de pages
        \setlength{\headheight}{14.2pt}% hauteur de l'en-tête
%%%%%%%%%%%%%%%%%%%style front%%%%%%%%%%%%%%%%%%%%%%%%%%%%%%%%%%%%%%%%% 
        \fancypagestyle{front}{%
                \fancyhf{}%on vide l'en-tête
                \fancyfoot[C]{\thepage}%
                \renewcommand{\headrulewidth}{0pt}%trait horizontal pour l'en-tête
                \renewcommand{\footrulewidth}{0pt}%trait horizontal pour le pied de page
                }
%\renewcommand{\chaptermark}[1]{%
%\markboth{\MakeUppercase{%
%\chaptername}\ \thechapter.%
%\ #1}{}}
%%%%%%%%%%%%%%%%%%%style main%%%%%%%%%%%%%%%%%%%%%%%%%%%%%%%%%%%%
        \fancypagestyle{main}{%
                \fancyhf{}
                \renewcommand{\headrulewidth}{0pt}%trait horizontal pour l'en-tête
                \renewcommand{\footrulewidth}{0pt}%trait horizontal pour le pied de pages
                \renewcommand{\chaptermark}[1]{\markboth{\chaptername\ \thechapter.\ ##1}{}}% redéfinition pour avoir ici les titres des chapitres des sections en minuscules
                \renewcommand{\sectionmark}[1]{\markright{\thesection\ ##1}}
                \fancyhead[C, RE, LO]{}
                \fancyhead[LE]{\rightmark}%
                \fancyhead[RO]{\leftmark}
                \fancyfoot[C]{}
                \fancyfoot[RO,LE]{\thepage}%
                \fancyfoot[LO,RE]{}
                }
%%%%%%%%%%%%%%%%%%%style back%%%%%%%%%%%%%%%%%%%%%%%%%%%%%%%%%%%%%%%%%  
        \fancypagestyle{back}{%
                \fancyhf{}%on vide l'en-tête
                \fancyfoot[C]{\thepage}%
                \renewcommand{\headrulewidth}{0pt}%trait horizontal pour l'en-tête
                \renewcommand{\footrulewidth}{0pt}%trait horizontal pour le pied de pages
                }

\newcommand{\gloss}[1]{{\sl \glssymbol{#1}}}
\newcommand{\Gloss}[1]{{\sl \Glssymbol{#1}}}

\newglossaryentry{mbd}
  {
	name=MBD,
	symbol=MBD,
	description=
	{
	  {\bf M}odel {\bf B}ased {\bf D}esign. Utilisé notamment dans des applications
	  automobiles, aérospatiales ou pour de l'informatique industrielle, MBD est un
	  méthode mathématique et visuelle pour comprendre et résoudre des problèmes
	   complexes de système de contrôle, de traitement du signal ou de système
	   communicant. Elle est beaucoup utilisée pour des logiciels embarqués.
	}
}
\newglossaryentry{rtw}
  {
	name=RTW-EC\up{\circledR},
	symbol=RTW-EC\up{\circledR},
	description=
	{
	  {\bf R}eal-{\bf T}ime {\bf W}orkshop-{\bf E}mbedded {\bf C}oder\up{\circledR}.
	  Aujourd'hui renommé Simulink Coder\up{\texttrademark}. Outil de génération de code
	  C et C++ à partir de diagrammes Simulink\up{\circledR}. Le code généré peut-être
	  utilisé pour des applications temps réel et hors temps réel, notamment pour le
	  prototypage rapide.
	}
}
\newglossaryentry{rtos}
  {
	name=OS temps réel,
	symbol=RTOS,
	description=
	{
	  ({\bf R}eal {\bf T}ime {\bf O}perating {\bf S}ystem en anglais). Système d'exploitation
	  multi-tâches destinée aux applications temps réel. Il facilite la création d'un système
	  temps réel via des ordonnances de tâches spécialisés afin de fournir aux développeurs
	  des systèmes temps réel les outils et les primitives nécessaires pour produire un
	  comportement temps réel souhaité dans le système final.
	}
}
\newglossaryentry{inca}
  {
	name=INCA,
	symbol=INCA,
	description=
	{
	  Logiciel de mesure et de calibration d'\gloss{ecu} via plusieurs protocoles de communication
	  développé par ETAS. Il permet de modifier les calibrations d'une \gloss{ecu} et de vérifier
	  que les changements au niveau des mesures sont corrects.
	}
}
\newglossaryentry{orianne}
  {
	name=Orianne,
	symbol=Orianne,
	description=
	{
	  {\bf O}util numé{\bf RI}que pourme m{\bf A}quettage de fo{\bf N}ctions de co{\bf N}trole mot{\bf E}eur.
	  Plate-forme de prototypage rapide de fonctions de contrôle moteur développée par Aboard Engineering. Voir \ref{sec:orianne}
	  pour plus de détails.
	}
}
\newglossaryentry{qgen}
  {
	name=QGen,
	symbol=QGen,
	description=
	{
	  Générateur open source de code embarqué temps réel en C et Ada à partir de modèles
	  Matlab\up{\textregistered} Simulink\up{\textregistered}. Basé sur
	  \gloss{projectp}, son développement est piloté par la société AdaCore.
	}
}
\newglossaryentry{projectp}
  {
	name=Project P,
	symbol=Project P,
	description=
	{
	  Projet open source visant à réduire les coût des approches de
	  développement \gloss{mbd} dans le cadre d'applications temps réel
	  embarquées. Le projet comprend de grands groupe industriel comme de plus
	  petites structure qui collaborent dans un but commun. Pour en savoir plus
	  {\tt http://www.open-do.org/projects/p/}
	}
}
\newglossaryentry{ecu}
  {
	name=ECU,
	symbol=ECU,
	description=
	{
	  {\bf E}ngine {\bf C}ontrol {\bf U}nit. Le boitier de contrôle moteur contenant des
	  pilotes et API pour manipuler les entrées sorties du boitier. L'application de contrôle
	  moteur sera déployé dans ce boitier et utilisera les pilotes de ce matériel pour interagir avec le moteur.
	}
}
\newglossaryentry{rte}
  {
	name=RTE,
	symbol=RTE,
	description=
	{
	  {\bf R}eal {\bf T}ime {\bf E}nvironment. Dans la plate-forme \gloss{orianne}, RTE
	  correspond à l'application développée à partir du code généré. En comparaison à
	  l'application gérant le matériel de l'\gloss{ecu} (pilotes, API, etc.).
	}
}
\newglossaryentry{can}
  {
	name=CAN,
	symbol=CAN,
	description=
	{
	  Le bus CAN ({\bf C}ontroller {\bf A}rea {\bf N}etwork) est un bus système série
	  très répandu dans beaucoup d'industries, notamment l'automobile. L'approche
	  bus permet de faire du multiplexage plutôt que de la communication point-à-point.
	}
}
\newglossaryentry{canalyzer}
  {
	name=CANalyzer,
	symbol=CANalyzer,
	description=
	{
	  Logiciel permettant de visualiser les trames émises et reçues sur un bus CAN développé par Vector. 
	}
}
\newglossaryentry{osgi}
  {
	name=OSGi,
	symbol=OSGi,
	description=
	{
	  {\bf O}pen {\bf S}ervices {\bf G}ateway {\bf i}nitiative est une plate-forme de services
	  fondée sur le langage Java qui permet de créer des composants dynamiques fournissant
	  des services. Ces composants peuvent se connecter et se déconnecter \og à la volée
	  \fg{} à la plate-forme.
	}
}
\newglossaryentry{maven}
  {
	name=Maven,
	symbol=Maven,
	description=
	{
	  Maven est un outil pour la gestion de l'automatisation de production des projets logiciels
	  Java. L'objectif recherché est de produire un logiciel à partir de se sources en optimisant
	  les tâches réalisées à cette fin et en garantissant le bon ordre de fabrication.
	}
}


\makeglossaries
